%% kapa: exam=1 round=25 year=2014 period=2 form=A subject=appraisal_law subject_ko=부동산관계법규 seq=4 number=4
%% source: Q-Net 감정평가사 기출문제 (cst00309) artlSeq=5101207 file=extracted/1008/5101207/2014년 제25회 감정평가사 1차시험 기출문제/2교시(A형).hwp
%% conversion: kordoc 4.24.51 -> tools/gen1.py (auto); review=proofread-sonnet; flags=none
\begin{question}{4}
국토의 계획 및 이용에 관한 법령상 광역도시계획에 관한 설명으로 옳은 것은?
\begin{choices}
  \item 특별시장ㆍ광역시장ㆍ특별자치시장ㆍ특별자치도지사ㆍ시장 또는 군수는 광역계획권을 지정할 수 있다.
  \item 광역계획권을 지정한 날부터 2년이 지날 때까지 관할 시ㆍ도지사로부터 광역도시계획의 승인 신청이 없는 경우에는 국토교통부장관이 수립한다.
  \item 광역계획권을 지정한 날부터 3년이 지날 때까지 관할 시장 또는 군수로부터 광역도시계획의 승인 신청이 없는 경우에는 국토교통부장관이 수립한다.
  \item 국가계획과 관련된 광역도시계획의 수립이 필요한 경우에는 국토교통부장관이 직접 또는 관계 중앙행정기관의 장과 공동으로 수립한다.
  \item 광역계획권이 둘 이상의 시ㆍ도의 관할 구역에 걸쳐 있는 경우에는 관할 시ㆍ도지사가 공동으로 수립한다.
\end{choices}
\end{question}
\answer{5}
